% ID: SEL-D0032-OA-E02
% ESTADO: cotejada. Segunda pasada visual y borrador compilado revisados.
% Fuente: Selectividad/matematicas_aplicadas/titular_junio_EXAMEN 5-ANDALUCÍA.pdf
% SHA-256: 5f24b200fb6cf34f6d72016ec3dc35787b4074de9b7752677e5b0d959ed0dfcd
% Página: 1; opción A; ejercicio original 2.

Sea $f(t)$ el porcentaje de ocupación de un determinado complejo hotelero en función del tiempo $t$, medido en meses, transcurrido desde su inauguración:
\[
f(t)=\begin{cases}
-\dfrac{5}{2}t^2+20t & \text{si }0\leq t\leq6\\[4pt]
\dfrac{90t-240}{t+4} & \text{si }t>6
\end{cases}
\]
\begin{enumerate}[label=\alph*)]
\item \textbf{(0.5 puntos)} \textquestiondown{}Evoluciona la función $f$ de forma continua?
\item \textbf{(0.5 puntos)} \textquestiondown{}Cuál sería el porcentaje de ocupación al finalizar el segundo año?
\item \textbf{(1 punto)} \textquestiondown{}En qué momentos el porcentaje de ocupación sería del 40\%?
\item \textbf{(0.5 puntos)} \textquestiondown{}Llegaría en algún momento a estar completo en caso de que estuviese abierto indefinidamente?
\end{enumerate}
